% Transcription — fonds Alexandre Grothendieck, shelfmark 10, batch 7 (pages 121-140).
% Established from the facsimile archives/batches/10/batch-07.pdf (a mirror of
% https://grothendieck.umontpellier.fr/10.pdf), per the skill
% transcribe-grothendieck.
%
% Pass: Opus 5 (claude-opus-5), 2026-09-12 — first pass, unchecked against the
% pages by a human.
%
% Pages 121-136 are the continuation, without a break, of the typescript
% « ⊗-catégories » begun page 117 in batch 6: the sentence cut off at the foot
% of page 120 (« Donc si k est ») resumes at the head of page 121. Sixteen
% leaves, typed, corrected throughout in his hand, running from the tail of
% item 7) to a « Corollaire » that closes the run. Page 137 is the blank verso
% of 136. Pages 138-139 are two leaves entirely in his hand, a fresh start
% headed « Motifs à coefficients dans un corps de nbs ». Page 140 is a brown
% folder cover inscribed by him, naming the run that begins after this batch.
%
%   121  end of 7): the graded category over a field, of finite type, gives
%        G_m, and C is recovered from Rep(G_m); then the same ⊗, φ, 1, u, v
%        with ψ changed for the Koszul rule (« Koko »), forced by the example
%        of H·(X,k) on compact spaces with Künneth; the claim, flagged « ??? »
%        in his margin, that C à la Koko then has no fibre functor over any
%        extension of k; 8) fibre functors valued in graded modules — « fonc-
%        teurs fibres tordus » — and the factorisation of F through Gradf(k)
%        = D^Z as a homomorphism of coalgebras A → k[T,T^{-1}], i.e. as
%        projectors π_i in U.
%   122  the relations π_i² = π_i, π_iπ_j = 0, Σπ_i = 1; a ⊗-structure on C
%        compatible with a), b) and with C → D equivalent to an associative
%        unital algebra structure on A for which A → k[T,T^{-1}] is a
%        homomorphism; the symmetry condition written as Δ(π_i) = Σ_{j+k=i}
%        π_j ⊗ π_k, and its reading: grading U by the degree of operators
%        makes the diagonal anticommutative, i.e. the multiplication of A
%        (of U) anticommutative; 9) the group of ⊗-automorphisms of the
%        forgetful functor is independent of a) to c) — the g ∈ U with
%        Δ(g) = g ⊗ g and augmentation 1.
%   123  the functor C → Rep_k(G) compatible with associative unital ⊗, and
%        when it is an equivalence; the data of G, of the grading G_m → G, and
%        the condition that the Koszul isomorphism M⊗N ⥲ N⊗M be compatible
%        with G; in car. 2 the distinction vanishes; it suffices that the
%        induced μ_2 → G be central, a fortiori that G_m → G be; the cancelled
%        hope that this is also sufficient; 10) the application to semi-simple
%        motives, the natural gradings on the H· coming from a grading of the
%        category, itself a translation of the algebraicity of the Künneth
%        correspondences — « On respire ! » — with the warning that these H·
%        are twisted fibre functors and not fibre functors.
%   124  for motives not necessarily semi-simple the category is filtered and
%        no longer graded; the wanted h: V → Grad(M) with the Koko structure,
%        and the wish that the cohomological functors factor through M as
%        genuine (untwisted) fibre functors, which is « pas idiot » since the
%        Koko symmetry was put on Grad(M) once and for all; or else replace
%        the geometric symmetry by the same one cancelled by (−1)^{ij};
%        11) intrinsic structure theorems in terms of gerbes.
%   125  « Voici quelques cogitations dans ce sens » — the symmetry data for a
%        given ⊗ on C with a faithful exact k-linear F, i.e. σ ∈ (U ⊗̂ U)^*
%        with σ'σ = 1 and σ(Δu)σ^{-1} = (Δu)' , the « twisted symmetry
%        condition on the diagonal map »; his box giving σ in the graded case
%        of n° 6 as π⁺⊗π⁺ + π⁺⊗π⁻ + π⁻⊗π⁺ − π⁻⊗π⁻; the unit condition
%        ε ∈ 1 + U⁺ ⊗̂ U⁺; and the associativity relation (id_U ⊗̂ Δ)(σ) =
%        σ^{12}σ^{13}.
%   126  base change to a k-algebra k', with U' and A'; the determination of
%        the exact k-linear functors G: C → Projf(k') with c: G(X⊗Y) ⥲
%        G(X)⊗G(Y), by a pro-module L = (L_α) over U' with L_α flat over U_α
%        and projective of finite type over k'; the isomorphism c as an
%        isomorphism of bi-U'-modules L ⊗̂ L ⥲ L ⊗̂ U'; B = Homcont_{k'}(L,k')
%        as a coalgebra over A'; his footnote (*) on the bimodule law, and his
%        margin asking that exactness = flatness be separated out first.
%   127  the structure as an isomorphism of comodules B ⊗_{k'} B ⥲ B ⊗_{k'} A'
%        and thence a composition law B ⊗_{k'} B → B compatible with the
%        A'-comodule structure; unit, associativity and commutativity of Δ
%        read off B being unital, associative and commutative; his margin, « NB
%        Ceci équivaut à dire que B → B ⊗_k A' est hom. d'alg. — A préciser »,
%        and « Ceci implique aussi que B est fidèlement plat sur k' »; in terms
%        of the affine monoid defined by A, everything is Spec B with the
%        monoid acting on the right.
%   128  the typed insertion at the head asking whether B is commutative when A
%        is, and whether the compatibility with symmetry is therefore
%        superfluous; the torsor picture over a faithfully flat k' → B; then A
%        merely associative unital with antipodism, σ ∈ U ⊗̂ U as above, the
%        category C' = C ⊗_k k' defined by the proring V = G(U) = L ⊗̂_U (U,int);
%        the question (1) whether a τ ∈ V ⊗̂_k V can be built from σ, the only
%        hope being that τ be invariant, U⁺·τ = 0; and what would follow —
%        that compatibilities with symmetry may be dropped in the study of
%        fibre functors, « celles-ci se récupérant toujours après coup ».
%   129  12. « Exemples instructifs » — G a finite monoid, U = k(G) with the
%        diagonal from the composition G×G → G, the Cartier-dual bialgebra to
%        the usual one; C as families (X_g) with Z_g = Σ_{g'g''=g} X_{g'} ⊗
%        Y_{g''} and the Dirac family as unit; the fibre functors given by
%        families (M_g) with M_{gg'} ≃ M_g ⊗ M_{g'}, whence rank 0 or 1;
%        a) a ⊗-functor exact and compatible with φ, ψ, 1 but not faithful,
%        possible only when G is not a group — G = (e,p) with p² = p, M_e = k,
%        M_p = 0; and his question at the foot whether a ⊗-functor compatible
%        with everything between ⊗-categories satisfying all the axioms,
%        rigidity included, is necessarily faithful and exact.
%   130  b) a ⊗-functor exact and faithful, compatible with 1, φ, ˅ but not
%        with ψ, not locally isomorphic to the given fibre functor: G a group,
%        the systems (M_g) equivalent to extensions of G by k^*, the functor
%        isomorphic to F iff the extension splits, the classes in bijection
%        with H²(G,k^*); even for G commutative and k algebraically closed H²
%        can be non-zero and a non-zero element dies in no extension of k;
%        hence fibre functors compatible with a given ψ and others not, and
%        the conclusion that the commutativity constraint cannot be neglected.
%   131  the commutativity constraints on this C correspond to the bilinear
%        antisymmetric functions G×G → k^*; when H²(G,k^*) is nul, every
%        non-zero σ gives a constraint admitting no fibre functor compatible
%        with it over any extension; graded fibre functors do not save it in
%        car. 2 and probably in none; the wanted example of a ⊗-category with
%        a commutativity constraint but none for which a fibre functor exists;
%        his note that this requires A commutative and U not.
%   132  c) X proper over k with H⁰(X,O_X) = k: Coh(X) has ⊗, φ, ψ, 1 and
%        internal Hom, but X ⊗ Y → Hom(X,Y) is not an isomorphism (rigidity
%        fails), so no true Galois theory; the fibre functors at the points of
%        X, defined over k(x), not exact and not faithful; the remedy by
%        coherent sheaves stratified on an open U, and the fact that for U = X
%        rigidity holds; « A prouver » — for C abelian with biadditive ⊗ and
%        rigidity, End(1) absolutely flat, the Hom of finite presentation, and
%        C an inductive limit of subcategories attached to composites of fields.
%   133  13) « Variante utile du tapis de n° 6 »: C not assumed additive, F not
%        assumed faithful, U = End_k(F), A the dual coalgebra, the canonical
%        factorisation through Modf(U/k) = Comodf(A/k); a bifunctor ⊗ with
%        F(X⊗Y) ≃ F(X)⊗F(Y) gives a diagonal on U; and the wanted converses.
%        The example named is the cohomology operations over F_p, where the
%        degrees in U are > 0 and Ũ = k + U⁺ is the canonical completion of a
%        graded ring with anticommutative diagonal.
%   134  « Théorème de ——— (rapporté par Husemöller) »: A a bialgebra over k,
%        associative and unital for both structures, cocomplete for its
%        canonical comodule filtration; M an A-module with a unital diagonal
%        M → M ⊗_k M that is A-linear, with a unit k → M a homomorphism of
%        unital coalgebras, M cocomplete, and u ↦ u.e injective; then M is a
%        free A-module, M ≃ A ⊗_k M̄ with M̄ = M ⊗_A k. Exemple 1: A
%        commutative, A the hyperalgebra of a connected formal group G, M the
%        distributions of a connected formal variety V, and the condition
%        saying that G acts on V with g ↦ g.a a monomorphism.
%   135  the action G×V → V×V a monomorphism by Nakayama, the quotient W = V/G
%        a G_W-torsor (SGA 3 VII_B), M̄ the coalgebra of W; in general the
%        torsor is not trivial and one has only that any lifting as k-vector
%        spaces gives A ⊗_k M̄ ≃ M; but G formally smooth gives a perfect
%        splitting, in particular in car. 0. Exemple 2: A commutative for its
%        algebra law, G an affine monoid scheme, the cocompleteness condition
%        read as a vanishing in n−1 of the variables, and the conjecture that
%        it amounts, for G of finite type, to G finite unipotent, and in
%        general to G prounipotent. The last third of the page is cancelled by
%        long diagonals.
%   136  Corollaire: if M → M ⊗_k M induces an isomorphism M ⊗_{(A,Δ_A)}
%        (A ⊗_k A) ⥲ M ⊗_k M, then u ↦ u.e is an isomorphism for all the
%        structures — the rank argument n = n², and the surjectivity of the
%        diagonal of M̄; and the conclusion that the local triviality
%        conjecture of his own notes on ⊗-categories holds for A finite over k
%        with cocomplete dual, provided B is not its own commutator ideal.
%   138-139  in his hand: C^{K'}, the objects of a ⊗-category over K carrying
%        an action of a finite extension K', with X ⊗_{K'} Y the quotient of
%        X ⊗_K Y and Hom_{K'} the corresponding kernel; when C comes from an
%        affine group G over K, C^{K'} comes from G_{K'} = G ⊗_K K'; then
%        motives with coefficients in a number field: E given by φ: K' →
%        End_Q(E_0), a polarisation giving Ě_0 ≃ E_0(ρ), and the boxed
%        (E_φ)^∨ = (E_{φ'})(ρ) — « ce n'est pas en général un isom.
%        (E_φ)^∨ ≃ E_φ(ρ) ! »; checked on level-zero, weight-zero K'-motives,
%        where rank one corresponds to continuous χ: π → K'^* and L_χ^∨ ≃ L_χ
%        exactly when χ takes values in ±1.
%
% Two material facts bear on the readings and are stated once here rather than
% page by page. The right-hand edge of the typescript is clipped in the
% digitisation: on the long-lined pages the last one or two characters of a
% line fall outside the image, and where a word is thereby cut the completion
% is given in \add{}. And several of his longer marginal notes are written
% slantwise down the left edge in a fast hand; their openings are legible and
% their continuations are not, and they are given as far as they go.
%
% The dating is the inventory's own, for the whole shelfmark.
\documentclass[11pt,a4paper]{article}
% Le préambule commun à toutes les transcriptions du fonds Grothendieck.
%
% Il tient en une idée : **l'incertitude se compose, elle ne se résout pas.**
% Une transcription qui lisse un mot douteux a détruit la seule chose pour
% laquelle on la faisait — un lecteur doit pouvoir savoir, à chaque endroit,
% ce qui était sur le feuillet et ce qui a été ajouté. Les six macros qui
% suivent sont donc l'appareil critique, et non de la décoration.
%
% Les mêmes commandes sont reconnues par `scripts/render.mjs`, qui produit la
% vue de lecture du site. Une macro ajoutée ici doit l'être aussi là-bas,
% faute de quoi le rendu échouera bruyamment — ce qui est voulu : un rendu
% silencieusement faux est pire qu'un rendu absent.

\ProvidesPackage{grothendieck}[2026/08/08 Transcriptions du fonds Grothendieck]

\usepackage{amsmath,amssymb,amsthm}
\usepackage{mathrsfs}
\usepackage{stmaryrd}
% Les diagrammes commutatifs sont partout dans ce fonds : c'est la langue
% même de la géométrie algébrique de Grothendieck, et une transcription qui
% les rendrait en prose ne transcrirait rien.
\usepackage{tikz-cd}
% Français par défaut — c'est la langue des feuillets — mais l'anglais est
% chargé aussi : la traduction et le résumé sont des documents anglais, et la
% typographie française (espaces avant les deux-points) y serait une faute.
% Ces documents basculent par \selectlanguage{english} en tête de corps.
\usepackage[english,french]{babel}
\usepackage{fontspec}
\usepackage{geometry}
\geometry{a4paper,margin=25mm,marginparwidth=28mm}
\usepackage{marginnote}
\usepackage{xcolor}
% ulem, not soul. `soul`'s \st and \ul are built on a character-by-character
% scan that XeLaTeX's font machinery breaks: the document fails with an
% undefined control sequence rather than degrading. ulem does the same two jobs
% and is Unicode-engine safe. `normalem` stops it hijacking \emph, which the
% transcriptions use for Grothendieck's own emphasis.
\usepackage[normalem]{ulem}
\usepackage{hyperref}
\hypersetup{colorlinks=true,linkcolor=black,urlcolor=black,pdfborder={0 0 0}}

% --- Métadonnées du lot -----------------------------------------------------
% Reprises telles quelles par le script de rendu, qui les affiche en tête de
% la vue de lecture. Elles fixent ce que le fichier prétend couvrir : sans
% elles, une transcription flotte sans point d'attache dans un fonds de seize
% mille feuillets.

\newcommand{\folder}[1]{\gdef\grothFolder{#1}}
\newcommand{\batch}[1]{\gdef\grothBatch{#1}}
\newcommand{\leaves}[2]{\gdef\grothFirst{#1}\gdef\grothLast{#2}}
\newcommand{\foldertitle}[1]{\gdef\grothFoldertitle{#1}}
\gdef\grothFolder{?}\gdef\grothBatch{?}\gdef\grothFirst{?}\gdef\grothLast{?}\gdef\grothFoldertitle{}

% --- Le résumé --------------------------------------------------------------
%
% Il ouvre la lecture modernisée, sous le titre « Résumé », et il est composé
% en petit. Ce n'est pas un ornement : c'est le seul passage du document qui
% ne soit pas des mathématiques, et un lecteur doit pouvoir le reconnaître
% avant de l'avoir lu — puis le sauter s'il connaît déjà le sujet, ou s'y
% arrêter s'il ne le connaît pas. Le filet qui le suit marque la reprise du
% corps.

\newenvironment{resume}
  {\small\setlength{\parskip}{0.45em}}
  {\par\medskip\hrule\medskip}

% --- Mots-clés ---------------------------------------------------------------
%
% En anglais, en clôture du résumé de la lecture modernisée : le vocabulaire
% moderne sous lequel chercher ce que les feuillets construisent. Le manifeste
% du site (`npm run manifest`) les extrait pour étiqueter la cote entière —
% cette ligne est la source unique des tags, il n'y a pas de fichier à côté.
\newcommand{\keywords}[1]{\par\smallskip\noindent
  {\footnotesize\textsc{Keywords} --- \textit{#1}}\par}

% --- L'appareil critique ----------------------------------------------------

% Le numéro de feuillet, en marge. C'est l'ancre qui permet de régler un
% désaccord sur une formule en regardant *un* feuillet plutôt qu'une chemise
% de six cents. Le site s'en sert aussi pour tourner les pages du fac-similé
% quand on fait défiler la transcription.
\newcommand{\leaf}[1]{%
  \marginnote{\footnotesize\color{gray!70}\textsf{#1}}%
  \label{leaf:#1}\ignorespaces}

% Illisible : on ne devine pas. Un mot inventé qui se lit comme les autres est
% le pire résultat possible.
\newcommand{\ill}{\textcolor{red!60!black}{[\,\ldots\,]}}

% Lecture douteuse : le mot est donné, et signalé comme incertain.
\newcommand{\uncertain}[1]{\uline{#1}}

% Ajout de l'éditeur — une lettre manquante, un mot sous-entendu.
\newcommand{\add}[1]{\textcolor{blue!50!black}{[#1]}}

% Biffé par Grothendieck lui-même. Ce qu'il a rayé fait partie du document :
% une rature dit souvent où la pensée a changé de direction.
\newcommand{\struck}[1]{\sout{#1}}

% Note du transcripteur, sur l'état matériel du feuillet.
\newcommand{\note}[1]{\par\smallskip{\small\color{gray!60!black}\textit{[#1]}}\par\smallskip}

% Note marginale de Grothendieck, distincte du corps du texte.
\newcommand{\marginal}[1]{\par\smallskip\noindent\hspace{1em}%
  {\small\color{gray!60!black}#1}\par\smallskip}

% --- Titre ------------------------------------------------------------------

\AtBeginDocument{%
  \begin{center}
    {\large\bfseries Fonds Alexandre Grothendieck --- cote n\textsuperscript{o}~\grothFolder}\\[2pt]
    {\itshape\grothFoldertitle}\\[4pt]
    {\small feuillets \grothFirst--\grothLast\ (lot \grothBatch)}\\[2pt]
    {\footnotesize\color{gray!60!black}Transcription établie d'après le fac-similé numérisé
     par l'Université de Montpellier.\\
     Les lectures incertaines sont soulignées, les illisibles notées [\,\ldots\,],
     les ajouts entre crochets.}
  \end{center}
  \vspace{1em}}

\endinput


\folder{10}
\batch{7}
\pages{121}{140}
\dating{[à partir de 1958]}
\watermark{Édition de démonstration}
\foldertitle{Catégories tensorielles : notes manuscrites (s.d.), tapuscrits (s.d.)}

\begin{document}

\section*{$\otimes$-catégories, suite (pages 121 à 136)}

\note{Suite sans rupture du tapuscrit commencé page 117 : la phrase laissée en
suspens au bas de la page 120, « Donc si $k$ est », se termine ici. Il est
corrigé de bout en bout à l'encre : ses insertions sont données entre crochets,
ses suppressions barrées, ses notes marginales en retrait. La machine n'avait
ni $\otimes$ ni lettres grecques ; on ne relève pas un à un les signes encrés
de sa main.}

\page{121}

corps et si on se borne aux modules de type fini sur $k$, on est dans les
conditions de 6), et on trouve le groupe algébrique $\underline{G}_{m}$ sur
$k$. D'ailleurs même sans condition sur $k$ et sans se borner à des modules de
type fini, on récupère $\underline{C}$ avec toutes ces structures par les
représentations de $\underline{G}_{m}$.

Gardons $\otimes$, $\varphi$ (et $1$, $u$, $v$) comme dessus, mais changeons
$\psi$ \struck{en lui} en utilisant la règle de Koszul. Ce \add{nouveau} choix
de $\psi$ s'impose dans certaines \add{questions}. Par exemple supposons $k$ un
corps, \struck{de car.\ nulle}, soit $\underline{V}$ la catégorie des espaces
topologiques compacts, avec le foncteur $\otimes$ donné pa\add{r le} produit
cartésien, et la donnée évidente d'associativité et de commutati\add{vité}
(pour le coup on n'a vraiment pas le choix !) ; il y a une unité, \add{à}
savoir l'espace ponctuel. Considérons le foncteur cohomologique
$X \mapsto H^{\bullet}(X,k)$ sur $\underline{V}^{\circ}$ ; utilisant Künneth on
trouve un isomorphisme $c$ de compatibilité aux $\otimes$ \add{de
$\underline{V}^{\circ}$ dans $\underline{C}$}. Ce dernier est compatible aux
structures a), b), c) à condition de choisir $\psi$ dans $\underline{C}$ par la
règle de Koko.\marginal{Cf.\ exemple 12}\note{« Koko » est son abréviation
plaisante pour Koszul ; elle revient aux pages 123 et 124.}

A prouver que lorsqu'on prend $\underline{C}$ comme dessus à la Koko, il n'y a
pas de foncteur fibre sur $\underline{C}$ sur aucune extension $K$ de $k$ (même
en se bornant aux vectoriels de dimension finie)
!\marginal{???}\marginal{déterminer tous les \ill{} qui sont compatibles
\ill{}}\note{Cette dernière note court en diagonale dans la marge inférieure
gauche ; son début se laisse lire, sa suite non.}

8) Pour pallier à cet inconvénient, il y a lieu \struck{lorsqu'on s'en va dans
une catégorie avec $\otimes$ : la catégorie des} d'étudier systématiquement les
$\otimes$-foncteurs compatibles avec a) b) c) à valeurs dans une catégorie
$\underline{C}$ \add{de modules gradués} comme ci-dessus : ce sont des
« foncteurs fibres tordus ». En particulier, si $\underline{C}$ est comme au
début de 6) (pas encore de $\otimes$), supposons que le foncteur $F$ se
factorise (de façon donnée) par les modules gradués, i.e.\ qu'on ait un
foncteur compatible avec les foncteurs oublis
$\underline{C} \to \operatorname{Gradf}(k) = \underline{D}^{\mathbb{Z}}$. Cela
correspond à un homomorphisme de coalgèbres de $A$ dans $k[T,T^{-1}]$, ou à la
donnée d'éléments $\pi_{i}$ de $U$ (les projections sur les composants
homogènes) satisfaisant les conditions habituelles

\page{122}

$\pi_{i}^{2} = \pi_{i}$, $\pi_{i}\pi_{j} = 0$ pour $i \neq j$,
$\sum_{i} \pi_{i} = 1$. La donnée d'une $\otimes$-structure sur $\underline{C}$
impliquant les données a) b), et celles-ci compatibles avec le foncteur
$\underline{C} \to \underline{D}$ (pour $c = \operatorname{id}$), équivaut alors
à une structure d'algèbre associative unitaire sur $A$, compatible avec la
structure de coalgèbre, et pour laquelle $A \to k[T,T^{-1}]$ est un
homomorphisme d'algèbres. Comment s'exprime le fait que \struck{\ill{}} sur
$\underline{C}$ il y a une structure de symétrie \struck{telle} respectée par le
foncteur $\underline{C} \to \underline{D}$ (et $c$ choisi plus haut) ? On peut
d'abord exprimer la structure \add{$\otimes$} précédente par une application
diagonale de $U$ satisfaisant
\[ \Delta(\pi_{i}) = \sum_{j+k=i} \pi_{j} \otimes \pi_{k} . \]
La condition cherchée revient à \struck{dire que si $M$ est un $U$-module de
degré $i$, i.e.\ tel que $\pi_{i}M = M$, et $N$ de degré $j$ i.e.\
$\pi_{j}N = N$, alors l'isomorphisme $x \otimes y \mapsto (-1)^{ij} y \otimes x$
de $M \otimes_{k} N$ sur $N \otimes_{k} M$ est compatible avec les opérations de
$U$} est celle-ci : si nous graduons $U$ par « degré des opérateurs dans des
gradués » i.e.\ \struck{\ill{}} ces quotients $U_{\alpha}$ par
$U_{\alpha}^{i} = \widehat{\sum} \pi_{i+j} U_{\alpha} \pi_{i}$, alors
l'application diagonale devient \emph{anticommutative} i.e.\ la multiplication
de $A$ \struck{\ill{}} est \emph{anticommutative}.\marginal{i.e.\ de $U$} (Il y
aurait lieu de prouver que la graduation de $A$ permet de récupérer la
\struck{\ill{}} graduation sur le foncteur oubli, i.e.\ les $\pi_{i} \in U$
i.e.\ $A \to k[T,T^{-1}]$, et que la graduation de $A$ \struck{\ill{}} peut
être choisie arbitrairement, avec la seule condition, de toutes façons
nécessaire, que $A$ soit anticommutative et l'homomorphisme diagonal
\struck{soit} compatible aux
graduations).\marginal{à vérifier (\ill{} pour \ill{}
?)}

9) Quelles que soient les structures a) \struck{à} c) en jeu,\marginal{(b)
— l'unité — étant \ill{}} le groupe des $\otimes$-automorphismes du foncteur
\struck{fibre} « oubli » en est indépendant, c'est le groupe des $g \in U$ tels
que $\Delta(g) = g \otimes g$ et d'augmentation $1$.\marginal{(et les hom.\
unitaires $g : A \to k$)} (NB cette dernière condition exprime la condition b)
de 3\textordmasculine). Remplaçant $k$ par une algèbre quelconque, on trouve
encore un \struck{\ill{}} schéma en monoïdes affine sur $k$, mais qui en général
ne permettra pas de reconstituer la situation entièrement. \struck{Donner
exemple} Il faut analyser de plus près

\page{123}

le foncteur compatible aux $\otimes$ associatifs unitaires qu'on obtient de
$\underline{C}$ dans $\operatorname{Rep}_{k}(G)$, et voir en particulier sous
quelles conditions il est une équivalence de catégories : \struck{Alors}
Lorsqu'il en est ainsi, la situation est décrite entièrement par la donnée de
$G$, de la graduation sur \struck{son} le foncteur oubli sur
\struck{$\underline{C}$} \add{$\operatorname{Rep}_{k}(G)$}, décrite par un
homomorphisme de monoïdes $\underline{G}_{m} \to G$, et enfin la condition que
pour deux représentations $M, N$ de $G$, l'isomorphisme de Koszul
$M \otimes N \xrightarrow{\sim} N \otimes M$ sur les vectoriels sous-jacents est
compatible avec l'action de $G$.\marginal{il \uncertain{f.\,d.\,p.} que $A$ soit
commutatif ! « géométriquement »}\marginal{dans le cas gradué anticommutatif}
Bien sûr, si $k$ est de car 2, ce n'est pas une restriction, et dans ce cas
d'ailleurs Koko ne change rien i.e.\ tous nos ennuis et distinguos
disparaissent. \struck{Si} Il suffit en tous cas que l'homomorphisme induit
$\mu_{2} \to G$ soit central ; à fortiori que $\underline{G}_{m} \to G$ soit
central.\marginal{et on voit que \struck{cela est} \ill{} si \ill{} implique
$A' = A^{\circ}$ \ill{}} \add{Cela signifie} \struck{i.e.} que la décomposition
d'un $M$ en composants \struck{homogènes} \add{pairs et impairs} soit invariante
sous $G$ i.e.\ se fasse dans $\underline{C}$, resp.\ qu'il en soit de même pour
la décomposition en degrés \add{quelconques}. Noter que cette dernière condition
s'exprime directement en termes de $\underline{C}$ (muni de sa structure
$\otimes$ et du foncteur \struck{oubli} gradué
$\underline{C} \to \operatorname{Grdf}(k)$), \struck{et on espère que cette
condition est suffisante (peut-être nécessaire et suffisante) pour qu'on soit
dans le cas favorable $\underline{C} \to \operatorname{Rep}_{k}(G)$ soit une
équivalence.}\marginal{si car.\ $\neq 2$}\note{Tout le bas de ce paragraphe, de
« et on espère » à « équivalence », est barré de plusieurs longues diagonales.}

10) Ce résultat s'appliquera alors, en particulier, à la catégorie des motifs
\add{semi-simples} (car les \struck{\ill{}} graduations naturelles sur les
foncteurs $H^{\bullet}$ qu'on arrive à y définir \struck{sont} proviennent d'une
graduation de la catégorie des motifs, --- ce qui n'est qu'une traduction du
fait que les correspondances de Künneth \add{cohomologiques} sur une variété
algébrique projective lisse sont algébriques). On respire ! Mais on fera
attention que les foncteurs $H^{\bullet}$ en question ne sont pas des foncteurs
fibres proprement, mais des foncteurs fibres tordus. (On ne pourra s'empêcher de
les appeler quand-même foncteurs fibres !) Quand on travaille avec des motifs
pas nécessairement

\page{124}

semi-simples, on a une catégorie qui n'est plus graduée, seulement filtrée. Si
$\underline{V}$ est la catégorie des schémas de type fini sur un corps fixé,
$\underline{M}$ la catégorie des motifs sur ce même corps (présumée avec toutes
les données et conditions de 1\textordmasculine), on doit avoir un foncteur
compatible avec ces structures (cohomologie motivique)
\[ h : \underline{V} \longrightarrow \operatorname{Grad}(\underline{M})
   \quad \text{(structure de Koko sur } \operatorname{Grad}(\underline{M})) , \]
et on désire que les foncteurs cohomologiques habituels sur $\underline{V}$ se
factorisent par $\underline{M}$. \struck{Ceci impose} \add{On voudra} que les
foncteurs \struck{fibres} obtenus sur $\underline{M}$ soient des foncteurs
fibres \struck{tordus seulement} véritables \add{(non tordus)}, ce qui n'est pas
idiot, \add{car} on avait mis la structure de symétrie de Koko sur
$\operatorname{Grad}(\underline{M})$ : on s'est débarassé ainsi une fois pour
toutes de la « torsion ».

On peut aussi, pour parler de vrais foncteurs fibres sur la catégorie des motifs
semi-simples, remplacer la donnée de symmétrie « géométrique » i.e.\ de Koko
pour le produit tensoriel dans \struck{$\operatorname{Grad}(\underline{M})$}
\add{la catégorie graduée $\underline{M}$}, par la même canulée par le signe
$(-1)^{ij}$ quand on a des arguments de degré $i, j$.

11) Bien sûr, il y a lieu de chercher des théorèmes de structure intrinsèques,
en termes de gerbes, pour des catégories $\underline{C}$ munies des structures
de 1\textordmasculine, \add{avec $k = \operatorname{End}(\underline{1})$ un
corps donné.} Il y a lieu au besoin de supposer qu'il existe une extension de
$k$ sur laquelle il y ait un foncteur fibre tordu gradué (ou multigradué ? Mais
où s'arrêter ?). Le plus beau serait évidemment de trouver un théorème de
structure général, sans \emph{supposer} l'existence d'un foncteur fibre plus ou
moins tordu où que ce soit, en définissant une notion assez générale de foncteur
fibre tordu (?) et en prouvant qu'il y en a toujours sur une extension
convenable de $k$.

\page{125}

Voici quelques cogitations dans ce sens, lorsque $\underline{C}$ est muni d'un
foncteur additif exact fidèle $k$-linéaire
$F : \underline{C} \to \operatorname{Modf}(k)$, et correspon\add{dant} à ce
titre à une algèbre profinie (ass.\ un.) $U$, resp.\ une coalgèbre ass.\ unit.\
$A$, sur le corps $k$ (cf.\ n\textordmasculine 6). On \struck{voit que la donnée
d'} suppose donné\add{e} une opération $\otimes$, compatible avec $F$, i.e.\
un\struck{e} hom.\ diagonal $U \to U \mathbin{\widehat{\otimes}} U$ i.e.\ une
composition $A \otimes_{k} A \to A$. On se propose de déterminer \add{toutes}
les données de symétrie $\psi$ pour $\otimes$ (\emph{pas seulement} celles
compatibles avec le foncteur $F$ et la donnée de symmétrie ordinaire sur
$\operatorname{Modf}(k)$ !) On trouve qu'elles correspondent aux
\[ (0) \qquad \sigma \in (U \mathbin{\widehat{\otimes}} U)^{*} \]
\marginal{Dans le cas du n\textordmasculine 6, \emph{gradué}, on aura
$\sigma = \pi^{+} \otimes \pi^{+} + \pi^{+} \otimes \pi^{-}
+ \pi^{-} \otimes \pi^{+} - \pi^{-} \otimes \pi^{-}$}
satisfaisant les conditions
\[ (1) \quad
\begin{cases}
\sigma'\sigma = 1 \ \text{(i.e. } \sigma \text{ et } \sigma' \text{ inverses
l'un de l'autre : on aura aussi } \sigma\sigma' = 1 , \\
\qquad \text{puisque } \sigma \text{ est supposé inversible)} \\[2pt]
\sigma(\Delta u)\sigma^{-1} = (\Delta u)' \quad \text{pour tout } u \in U
\end{cases} \]
\marginal{condition de symétrie « tordue » de l'application diagonale}
en désignant par $x \mapsto x'$ l'automorphisme de symétrie de
$U \mathbin{\widehat{\otimes}} U$. NB On ne dit pas que $\Delta U$ est stable
sous $x \mapsto x'$, mais on considère $\Delta$ et le symétrique $\Delta'$
défini par $\Delta'(u) = \Delta(u)'$, qui définit la loi $\otimes$
\emph{opposée} de celle définie par $\Delta$, et les relations entre $\sigma$ et
$\sigma'$ sont symétriques et s'expriment \add{en disant} que $\sigma$
transforme $\Delta$ en $\Delta'$, $\sigma'$ transforme $\Delta'$ en $\Delta$.

Lorsqu'il y a une unité bilatère pour $\otimes$ et que $F$ est compatible avec
celle-ci, \add{i.e.\ $A$ admet une unité bilatère,} \struck{alors} la
compatibilité de $\psi$ à l'unité de $\underline{C}$ s'exprime par les relations
\[ (2) \qquad (\operatorname{id}_{U} \otimes \varepsilon_{U})(\sigma)
 = (\varepsilon_{U} \otimes \operatorname{id}_{U})(\sigma) = 1 ,
 \quad \text{i.e.} \quad
 \varepsilon \in 1 + U^{+} \mathbin{\widehat{\otimes}} U^{+} . \]

Lorsqu'on a une donnée d'associativité $\varphi$ pour $\otimes$, compatible avec
$F$ et l'associativité ordinaire de $\otimes$ sur $\operatorname{Modf}(k)$,
i.e.\ $A$ est associative, alors la compatibilité entre $\varphi$ et $\psi$
s'exprime par une relation du style
\[ (3) \qquad \struck{\ill{}} \quad
 (\operatorname{id}_{U} \mathbin{\widehat{\otimes}} \Delta)(\sigma)
 \struck{\ill{}} = \sigma^{12}\sigma^{13} \]
\marginal{où les $\sigma^{ij}$ sont les 3 images de $\sigma$ dans
$U \otimes U \otimes U$ par $x \otimes y \mapsto x \otimes y \otimes 1$ et les
deux autres types analogues.}

\page{126}

\struck{où les $\sigma_{i}$ (resp.\ $\sigma'_{i}$) sont déduits de $\sigma$
(resp.\ de $\sigma'$) par les trois homomorphismes
$\delta_{i} : U \mathbin{\widehat{\otimes}}_{k} U \to
U \mathbin{\widehat{\otimes}}_{k} U \mathbin{\widehat{\otimes}}_{k} U$
\struck{utilisant} $\Delta \otimes \operatorname{id}_{U}$ et symétriques.} (NB je
n'ai pas écrit le calcul et ne garantis pas la forme exacte de la relation de
l'hexagone).

Soit $k'$ une algèbre (comm.\ ass.\ unitaire !) sur $k$, d'où
$U' = (U'_{\alpha})$ et $A' = \varinjlim_{\alpha} A'_{\alpha}$ sur $k'$, avec
leurs structures \struck{d'algèbre} bigèbres (resp.\ bi-progèbre pour $U'$). On
se propose de déterminer les foncteurs
\[ G : \underline{C} \to \struck{\operatorname{Modf}}
   \operatorname{Projf}(k') \]
\marginal{(modules proj de t\,f sur $k'$)} \emph{exacts et $k$-linéaires}, munis
d'un isomorphisme $c : G(X \otimes Y) \xrightarrow{\sim} G(X) \otimes_{k} G(Y)$.
\struck{\ill{}} Tout d'abord la donnée de $G$ exact à droite
\marginal{$\underline{C} \to \operatorname{Mod}(k')$} revient à la donnée d'un
$L = (L_{\alpha})$, module \add{à droite} sur $U$, muni d'une structure de
$k'$-module \struck{\ill{}} commutant aux opérations de $U$ (NB
$L_{\alpha} = G(U_{\alpha})$), par
\[ G(M) = L \otimes_{U} M . \]

La $k$-linéarité signifie que les deux structures sur $k$ de $L$ coïncident,
i.e.\ que $L$ provient d'un module (pro) sur $U' = (U'_{\alpha})$ (chaque
$L_{\alpha}$ est un module sur $U'_{\alpha}$). \add{On peut alors écrire
$G(X) = L \otimes_{U'} X'$, où $X' = X \otimes_{k} k'$.} \struck{Le fait que
l'on s'envoie dans $\operatorname{Projf}(k')$ s'exprime} \add{et que le foncteur
envisagé $G$ est exact} s'exprime alors par le fait que les $L_{\alpha}$ sont
\struck{libres} \add{proj.} de type fini \emph{sur $k'$}, et que les
$L_{\alpha}$ sont \emph{plats sur les $U_{\alpha}$} (nécessité triviale,
suffisance un petit argument par réduction au cas $X$ simple, puis $U$ un corps
égal à $X$ …). Donc on trouve exactement les \emph{modules
$L = (L_{\alpha})$ sur $U'$ tels que les $L_{\alpha}$ soient plats sur
$U_{\alpha}$ et proj de tf sur} $k'$.\marginal{séparer les conditions, commencer
par exactitude $=$ platitude sur les $U_{\alpha}$} D'autre part, indépendamment
de ces conditions de platitude, la donnée d'un isomorphisme $c$ revient à la
donnée d'un isomorphisme de bi-$U'$-modules \add{à droite}
\marginal{si dans $\underline{C}$ il y a $1$ et $^{\vee}$ compatibles \ill{}
$F$, i.e.\ $\Delta u$ unitaire et \struck{\ill{}} antipodisme}
\[ L \mathbin{\widehat{\otimes}}_{\uncertain{k'}} L \xleftarrow{\ \sim\ }
   L \mathbin{\widehat{\otimes}}_{k'} U'
   \quad \bigl( \simeq L \otimes_{(U,\Delta)}
   (U \mathbin{\widehat{\otimes}}_{k} U)
   \xrightarrow{\ \sim\ } L \otimes_{(U',\Delta')}
   (U' \mathbin{\widehat{\otimes}}_{k} U') \bigr) . \]
\marginal{(grâce à l'antipodisme)
$x \otimes u \mapsto x \otimes (u \otimes 1)$}

\struck{En termes} Utilisant les conditions de platitude, on peut définir
$B = \operatorname{Homcont}_{k'}(L,k')$, qui est limite inductive des
$B_{\alpha}$ proj de tf sur $k'$ (à fortiori \add{il est} plat sur $k'$), et une
coalgèbre sur $A'$ (limite de coalgèbres

\note{Une note de sa main occupe le bas de la page : « ($*$) La loi de bim.\ sur
$L \otimes_{U} U'$ est donnée par \{ 1\textsuperscript{ère} loi : déduite de
celle de \uncertain{$U'$} ; 2\textsuperscript{ème} loi : donnée par celle du
mod.\ : d'abord sur $U \otimes U'^{\circ}$ [$U'^{\circ}$ opère à droite sur
$U'$ !] via $(\operatorname{id}_{U} \otimes g)(\Delta u)$ \}. » Une seconde note
court en diagonale dans la marge gauche et ne se laisse pas lire.}

\page{127}

$B_{\alpha}$ sur les $A'_{\alpha}$ qui se réduisent les unes en les autres
…), \struck{les $B$} de plus chaque $B_{\alpha}$ étant coplat sur
$A_{\alpha}$. La structure précédente revient à un isomorphisme de
\add{co}modules sur \struck{\ill{}} $A'$
\[ (*) \qquad B \otimes_{k'} B \xleftarrow{\ \sim\ } B \otimes_{k'} A' . \]
\struck{En} Utilisant l'augmentation du facteur $A'$, on en conclut un
homomorphisme
\[ (**) \qquad B \otimes_{k'} B \longrightarrow B \]
qui sera un \emph{homomorphisme de comodules} sur $A'$ (NB utiliser la
multiplication de $A'$ pour définir la structure de $A'$-comodule sur le produit
tensoriel sur $k$ de deux tels), i.e.\ on a une loi de composition sur $B$
« compatible avec la structure de $A'$-comodule ».\marginal{Sauf erreur, la
donnée de $(*)$ \emph{équivaut à celle de} $(**)$}\marginal{NB Ceci
équivaut à dire que $B \to B \otimes_{k} A'$ est hom.\ d'alg.}\marginal{A
préciser}

Si $\Delta$ sur $U$ est unitaire i.e.\ $A$ est unitaire i.e.\ le $\otimes$ sur
$X$ \struck{\ill{}} admet une unité bilatère compatible avec $F$, alors la
compatibilité de $c$ avec les unités s'exprime, sauf erreur, par la condition
que la loi de composition $(**)$ est unitaire. En effet, il faut écrire que l'on
a un isomorphisme \struck{\ill{}} $L \otimes_{U} k' \xrightarrow{\sim} k'$
satisfaisant certaines conditions, or cela revient à se donner une
\struck{élément de} augmentation de $L$ satisfaisant une certaine condition
(savoir \struck{être compatible avec} \add{être unité bilatère pour} la
structure diagonale de $L$) ou encore un élément de $B$ qui est une unité
bilatère.\marginal{Ceci implique aussi que $B$ est \emph{fidèlement} plat
sur $k'$.} De même, si $\Delta$ est associatif, i.e.\ \add{si} on a pour
$\otimes$ une donnée d'associatilité $\varphi$ compatible avec $F$, alors $c$
est compatible avec celle-ci sss $B$ est une algèbre associative. \struck{Quand
les deux conditions sont vérifiées} Enfin, quand $\Delta$ est commutative i.e.\
$\underline{C}$ avec $\psi$ donnée de symétrie, …{} alors $c$ est compatible
avec cette dernière, et la donnée de symétrie habituelle sur
$\operatorname{Projf}(k')$ sss la loi de composition de $B$ est
\struck{\ill{}} commutative. (Mais alors, en termes du monoïde affine $H$
défini par $A$ sur $k$, on voit que les structures de $B$ s'expriment en termes
de $\operatorname{Spec} B$ en disant que $H_{k'}$ y opère à droite, et on voit de
plus qu'après le changement

\page{128}

\note{Le haut du feuillet porte, tapé en interligne simple et précédé d'un
« NB » de sa main, l'alinéa suivant, inséré avant la reprise de la phrase de la
page 127.}

\emph{Question} Sous réserve $B$ associative unitaire, si $A$ est commutative,
$B$ l'est-elle aussi ? (On suppose $A$ associative unitaire). Question analogue
quand on a des structures graduées avec anticommutativité. Donc, la condition de
\struck{symé}comp.\ aux sym.\ pour un $c$ est-elle surabondante, en prés.\ des
autres compat.\ ?

de base fidèlement plat $k' \to B$, \struck{\ill{}} il devient trivial, donc on
a un « torseur » sous le monoïde $H_{k'}$ ; alors le nouveau foncteur fibre $G$
est déduit de $F' = F \otimes_{k} k'$ par torsion à l'aide de $G$ …

Mais supposons seulement $A$ algèbre associative unitaire, et $B$ itou, sans
condition de commutativité. \add{On suppose $A$ avec \emph{antipodisme}
(s'en passer ?).} Donnons-nous un
$\sigma \in U \mathbin{\widehat{\otimes}}_{k} U$ \struck{\ill{}} comme
ci-dessus pour introduire une relation de commutativité $\psi$. Notons (qu'à
celà près que $k'$ n'est plus supposé un corps) on trouve une catégorie
$\underline{C}' = \underline{C} \otimes_{k} k'$ sur $k'$, avec $\otimes$,
définie par le « proanneau »\marginal{\ill{} $U'$, \ill{}}
\[ V = G(U) = L \mathbin{\widehat{\otimes}}_{U} (U,\operatorname{int}) , \]
où $(U,\operatorname{int})$ désigne $U$ muni de la structure de $U$-module
\struck{\ill{}} (et même d'objet « pro-anneau » (pour $\otimes$) de
$\underline{C}$) défini à l'aide de sa structure de bimodule sur
$U^{\circ} \mathbin{\widehat{\otimes}} U$ \struck{\ill{}} par restriction des
scalaires via
$U \xrightarrow{\Delta} U \mathbin{\widehat{\otimes}} U
\xrightarrow{g \otimes \operatorname{id}_{U}}
U^{\circ} \mathbin{\widehat{\otimes}} U$ : Tordant par $L$ et utilisant le fait
que l'on a une compatibilité avec $\otimes$, on trouve un pro-anneau $V$ sur
$k'$, et ce dernier hérite de $U$ une application diagonale …

Ceci dit, si on a une donnée
$\sigma \in U \mathbin{\widehat{\otimes}} U$ satisfaisant les conditions
\add{plus haut}, (1) peut-on définir en termes de $\sigma$ (ou existe-t-il tout
au moins) un élément $\tau \in V \mathbin{\widehat{\otimes}}_{k} V$ satisfaisant
les relations analogues à (1), de telle façon que $(G,c)$ soit compatible aux
isomorphismes de symétrie $\psi$ dans $\underline{C}$ défini par $\sigma$,
$\psi'$ dans $\underline{C}'$ défini par $\tau$ ??? La seule chance de
\struck{les montrer} construire un $\tau$ canoniquement semble de prouver que
$\tau$ est « invariant » lorsqu'on fait opérer $U$ sur
$U \mathbin{\widehat{\otimes}}_{k} U = (U,\operatorname{int})
\mathbin{\widehat{\otimes}} (U,\operatorname{int})$ grâce à son \struck{\ill{}}
action dite plus haut sur les deux facteurs, ce qui s'exprime par la relation
$U^{+}\!\cdot\tau = 0$.\marginal{Vérifier dans le cas gradué ce qui en est !} ---
Si ceci marchait (ce qui améliorerait le résultat souhaité en NB au haut de la
page), cela montrerait que dans l'étude des foncteurs fibres, il y a lieu de
laisser tomber les questions de compatibilité aux symétries, celles-ci se
récupérant toujours après coup.\marginal{Mais, cf.\ exemples n\textordmasculine
11}

Il resterait alors la question d'une description \emph{intrinsèque} d'une
catégorie $\underline{C}$ comme dans le n\textordmasculine 1, \add{avec $X$}
indépendamment de la donnée d'un $F$ préli-

\note{Une note de sa main, en trois alinéas a), b), c), court en diagonale sur
toute la marge gauche. Seuls les débuts se laissent lire : « a) donner une
\ill{} intrinsèque : $\underline{C} \otimes_{k} k'$, indépendante \ill{} de $F$
(NB $k$, $k'$ \ill{} une catégorie abélienne \uncertain{additive}) » ;
« b) Montrer \ill{} $\otimes$ \ill{} » ;
« c) $G : \underline{C} \to \operatorname{Mod}(k')$ \ill{} ».}

\page{129}

minaire, --- sous réserve tout au plus qu'il existe un tel $F$ sur une extension
convenable de $k$ (chose peut-être toujours vraie).

\subsection*{12. Exemples instructifs.}

Soit $G$ un monoïde fini. Contrairement au yoga galoisien général, au lieu de
prendre $A = k(G)$ (anneau des fonctions sur l'ensemble sous-jacent à $G$)
prenons $U = k(G)$, avec l'application diagonale déduite de
$G \times G \to G$ (la composition). (On a la bialgèbre « duale » \add{au sens
Cartier} de celle qu'on considère d'habitude). Donc la catégorie
$\underline{C}$ définie par $G$ peut être considérée comme la catégorie des
familles $(X_{g})_{g \in G}$ de vectoriels indexés par $G$, le produit
\add{tens.} \struck{l'application diagonale} de deux tels $X$ et $Y$ étant $Z$
défini par
\[ Z_{g} = \sum_{g'g''=g} X_{g'} \otimes Y_{g''} . \]
L'unité est la « famille de Dirac » définie par $X_{g} = 0$ si $g \neq e$,
$X_{e} = k$. On a une catégorie avec $\otimes$, $\varphi$, $1$, et de plus un
$\psi$ si $G$ est commutatif, un $^{\vee}$ si $G$ est \add{un groupe}. Un $L$
avec une application diagonale $L \to L \otimes_{k} L$ compatible avec
$U \to U \mathbin{\widehat{\otimes}}_{k} U$ et donnant un isomorphisme par
extensions des anneaux de base, peut s'interpréter comme une famille
$(M_{g})_{g \in G}$, avec des isomorphismes donnés
\[ (*) \qquad M_{gg'} \simeq M_{g} \otimes M_{g'} . \]
\marginal{Ceci implique que le rang des $M_{g}$ est 0 ou 1}
C'est compatible avec $1$ si on a un isomorphisme $M_{e} \simeq k$
\add{$(**)$}, qui rende les isomorphismes $(*)$ triviaux quand
\struck{\ill{}} $g$ ou $g'$ est $= 1$. C'est associatif si et seulement si les
isomorphismes $(*)$ satisfont la \struck{condition} conditions
d'associativité évidente. De même, $G$ étant commutatif, pour la commutativité.
--- Le foncteur $F'$ est toujours exact, car $M$ est plat sur $U$.

a) Exemple d'un $\otimes$-foncteur exact, compatible avec $\varphi$, $\psi$,
$1$, \emph{mais non fidèle} (donc pas loc.\ isom.\ au $\otimes$-foncteur
donné). Ce n'est possible dans notre exemple que si $G$ n'est pas un groupe
(\struck{sinon $M$} car $M_{g}$ est de rang 1 pour $g$ inversible). Mais prenons
$G = (e,p)$, avec $p^{2} = p$, prenons $M_{e} = k$, $M_{p} = 0$. On n'a plus le
choix pour les isomorphismes $(*)$, qui satisfont toutes les compatibilités
(avec $1$ par construction, $\varphi$, $\psi$).

\note{Au bas de la page, de sa main : « \emph{Question} Un $\otimes$-foncteur
compatible avec tout entre $\otimes$-catégories satisfaisant à tous les axiomes
(y compris rigidité) est-il néc.\ fidèle et exact ?? »}

\page{130}

b) Exemple d'un $\otimes$ foncteur exact fidèle, compatible avec
\struck{$\varphi$} $1$, $\varphi$, $^{\vee}$ (mais pas \struck{avec $\psi$} !),
et qui n'est pas localement isomorphe au foncteur fibre donné. --- On prend pour
$G$ un groupe. On note alors que la donnée d'un système de $M_{g}$ et
d'isomorphismes $(*)$ \struck{\ill{}} $(**)$, satisfaisant les diverses
compatibilités, est équivalente à la donnée d'une extension de $G$ par $k^{*}$.
Le foncteur \add{$F'$} est isomorphe au foncteur $F$ sss cette extension
splitte. L'ensemble des classes d'isomorphie de foncteurs fibres définis sur $k$
est ainsi en correspondance 1--1 avec $H^{2}(G,k^{*})$. Or Même si $G$ est
commutatif (par exemple le produit de deux groupes cycliques isomorphes d'ordre
premier à \struck{\ill{}} car.$k$) et $k$ algébriquement clos, le $H^{2}$ peut
être non nul ; \struck{\ill{}} d'ailleurs un élément non nul de $H^{2}$ ne se
tue dans aucune extension de $k$ (si $k$ alg.\ clos), d'après le principe de
l'extension finie.

Si cependant $G$ est commutatif, et si on se borne aux $F'$ compatibles avec
$\psi$ également, on trouve les extensions commutatives de $G$ par $k^{*}$ ; si
$k$ est algébriquement clos, donc $k^{*}$ divisible, alors toute telle extension
splitte, donc $F'$ est alors isomorphe à $F$ (en conformité avec la théorie
galoisienne relative au cas où l'application diagonale de $U$ est commutative
…).

Ainsi, il peut y avoir des foncteurs fibres compatibles avec $\psi$ donné, et
d'autres qui ne le sont pas et qui par suite ne sont pas loc.\ isom.\ au
premier. Cela montre donc qu'on ne peut négliger la \struck{condition de} la
contrainte de commutativité comme on aurait aimé pouvoir le faire. De plus, cela
suggère qu'en l'absence de contrainte de commutativité, la classification des
foncteurs fibres se fait par des groupes du type de cohomologie de Brauer.

\note{Dans la marge gauche, en diagonale : « étudier les contraintes \ill{}
associées \ill{} mod 2 », puis, encadré, « Exemple de foncteurs fibres
\uncertain{gradués} mod 2 non loc.\ isom.\ aux \ill{} ?? ».}

\page{131}

Dans cet exemple, on voit tout de suite qu'il y a sur $\underline{C}$ une
contrainte de commutativité compatible avec les autres données sss $U$ est
commutatif c'est à dire $G$ est commutatif, et que lesdites contraintes
correspondent \add{alors} aux fonctions $G \times G \to k^{*}$ qui sont
bilinéaires et antisymétriques. Dans le cas où le $H^{2}(G,k^{*})$ est nul (par
exemple $G$ cyclique \add{géométriquement} \struck{et $k$ algébriquement clos})
on voit donc que pour tout choix non nul de $\sigma$, on a une contrainte de
commutativité pour laquelle il n'y a aucun foncteur fibre compatible avec ladite
contrainte (et les autres) sur aucune extension de $k$. On peut songer à s'en
tirer en trouvant au moins des foncteurs fibres gradués compatibles à toutes les
contraintes, mais on voit déjà que c'est impossible si $k$ de car.\ deux (où il
n'y a pas de différence entre l'existence de l'un et l'autre type de foncteur
fibre, gradué ou non) ; et sans doute l'analyse doit montrer que l'on a un même
résultat négatif en toute caractéristique. Noter cependant dans ces exemples on
pouvait au moins, en modifiant la contrainte de commutativité donnée, en trouver
une qui admette un foncteur fibre compatible à ladite contrainte. Il y aurait
donc lieu de donner un exemple d'une $\otimes$-catégorie définissable par une
bialgèbre, admettant une contrainte de commutativité (compatible au reste) mais
aucune contrainte de commutativité pour laquelle on puisse trouver un foncteur
fibre (gradué au besoin) sur aucune extension du corps donné.

\note{De sa main, au bas de la page : « Avec les notations du
n\textordmasculine 11, il faudra donc $\Delta_{U}$ [i.e.\ $A$] commutatif, et
$U$ non commutatif (sinon l'existence de $\sigma$ impliquerait la commutativité
de $\Delta$) ]. »}

\page{132}

c) Soit $X$ un schéma \struck{complet} \add{propre} sur un corps $k$, supposons
$H^{0}(X,\underline{O}_{X}) = k$. La catégorie des faisceaux cohérents sur $X$
est une catégorie avec $\otimes$, $\varphi$, $\psi$, $1$,
$\underline{\operatorname{Hom}}$ internes, mais l'homomorphisme
$\check{X} \otimes Y \to \underline{\operatorname{Hom}}(X,Y)$
\emph{n'est pas un isomorphisme}.\marginal{$\operatorname{Hom}(X,Y)$ de
dim.\ finie sur $k$}\marginal{(axiome de rigidité)} C'est pour cela qu'il n'y a
pas lieu de s'attendre à une vraie théorie de Galois de $\underline{C}$. Il y a
des foncteurs fibres associés aux divers points de $X$ (pas nécessairement
fermés), à $x$ correspond un foncteur fibre défini sur $k(x)$ ; sans doute il
n'y en a pas d'autres à isomorphisme près (à vérifier). Ils ont toutes les
compatibilités, mais \emph{ils ne sont pas exacts}, \struck{\ill{}}
\emph{ni fidèles}. Pour rémédier au défaut d'exactitude, sans que
$\underline{C}$ cesse d'être abélienne (sinon on prendrait les Modules loc.\
libres seulement !), on peut se fixer un ouvert $U \subset X$ et définir
$\underline{C}$ en prenant des Faisceaux cohérents qui, sur $U$, sont munis d'une
stratification relativement à $k$ : tout point $x \in U$ fournit un foncteur
fibre exact. Cependant, si $U \neq X$, ce foncteur ne sera pas fidèle. Bien sûr,
si $U = X$, l'axiome de rigidité est satisfait, et on sait qu'on a une théorie
de Galois. Il doit résulter de cette dernière que $\underline{C}$ n'a pas de
foncteur fibre (compatible avec toutes les données) sur une extension de $k$
(car autrement l'axiome de rigidité serait satisfait).

\emph{A prouver} : Soit $\underline{C}$ une catégorie
\add{abélienne} avec $\otimes$ biadditif, $\varphi$, $\psi$, $1$, $^{\vee}$ et
rigidité, \add{$k = \operatorname{End}(\underline{1})$}. Alors
a) $\operatorname{End}(\underline{1})$ est un anneau absolument plat, donc un
composé fini de \struck{groupes} corps s'il est noethérien. b) Les
$\operatorname{Hom}(X,Y)$ sont des modules de présentation finie sur $k$
\marginal{i.e.\ cohérents} (NB les modules de présentation finie sur $k$ forment
une catégorie abélienne, en fait $k$ est un anneau « cohérent »).
c) \struck{Toute sous-catégorie de $\underline{C}$ stable par tout et à
engendrement fini} $\underline{C}$ est limite inductive de sous-catégories (en
général pas pleines) stables par tout, y compris noyaux, conoyaux, $^{\vee}$,
$\otimes$, et correspondants à des anneaux $k_{i}$ qui sont des
\struck{sommes} composés de corps. d) Le cas d'un composé de corps se ramène à
celui d'un corps.

\page{133}

13) \emph{Variante utile du tapis de n\textordmasculine 6.}\note{Le chiffre tapé
est 12 ; il est corrigé en 13 de sa main.} Soit $\underline{C}$ une catégorie,
et \struck{\ill{}}
\[ F : \underline{C} \longrightarrow \operatorname{Modf}(k) \]
un foncteur, où $k$ est un corps. On ne suppose pas $\underline{C}$ additive, ni
$F$ fidèle. On définit alors l'anneau profini
$U = \operatorname{End}_{k}(F)$, la coalgèbre duale $A$, et on trouve une
factorisation canonique de $F$ à travers
$\operatorname{Modf}(U/k) = \operatorname{Comodf}(A/k)$. Supposons donné dans
$\underline{C}$ un bifoncteur $\otimes$, \struck{alors} \add{et} un isomorphisme
de bifoncteurs $F(X \otimes Y) \simeq F(X) \otimes F(Y)$, alors on trouve une
application diagonale sur $U$ i.e.\ une loi de composition \struck{\ill{}} sur
$A$ : en effet, on constate que l'anneau des endomorphismes du bifoncteur du
premier membre est $U \mathbin{\widehat{\otimes}}_{k} U$, et $U$ s'envoie dedans
grâce à l'isomorphisme donné $F(X \otimes Y) \simeq F(X) \otimes F(Y)$. Si on se
donne pour $\otimes$ des données d'unité, d'associativité ou de commutativité,
la compatibilité de ces données avec $c$ s'exprime alors comme il est dit dans
n\textordmasculine 6,\marginal{en termes de $\Delta$} et il y a lieu de donner
des variantes graduées (si possible plus générales et astucieuses qu'au
n\textordmasculine 8). En fait, les résultats spécifiques au
n\textordmasculine 6 seraient plutôt, dans le cas d'un foncteur fidèle (et
additif \struck{\ill{}} exact de catégories abéliennes ? à voir de quoi
exactement on a besoin), des réciproques aux constructions précédentes, i.e.\
comment déduire d'une application diagonale sur $U$ un foncteur $\otimes$ sur
$\underline{C}$, dont les propriétés vont réfléter celles de cette dernière.

\emph{Exemple intéressant} : les opérations cohomologiques sur
$\underline{F}_{p}$. Il mérite quelques explications : sauf le morceau engendré
par les $\pi_{i}$ \struck{\ill{}} les degrés dans $U$ sont $> 0$, et
$\widetilde{U}$\marginal{$= k + U^{+}$} est le complété canonique d'un anneau
gradué à application diagonale anticommutative, donc $\widetilde{A}$ est un
anneau analogue. Dire comment on récupère $U, A$ à partir de $\widetilde{U}$,
$\widetilde{A}$.

\page{134}

\subsection*{Théorème de \ill{} (rapporté par Husemöller).}

\note{Le tapuscrit laisse le nom en blanc --- une suite de tirets --- et ne le
complète pas.}

Soit $A$ une bialgèbre sur le corps $k$, $A$ associative et unitaire pour ses
deux structures, « cocomplète » pour sa filtration canonique de comodule (duale
de la filtration par les puissances de l'idéal d'augmentation du dual ; de façon
précise, écrivant $\otimes^{n} A = \otimes^{n} A^{+} + $ fourbis, d'où via la
m.ème itérée du morphisme diagonal, une application
$A \to \otimes^{n} I^{+}$, dont le noyau est
$\operatorname{Filt}_{m}(A)$).\marginal{intersection avec $A^{+}$} (NB
l'hypothèse cocomplète est satisfaite pour une coalgèbre à degrés $\geq 0$
\add{avec unité} réduite à $k$ en degré $0$, associative et unitaire, ---
parait-il !). Soit $M$ un $A$-module, muni d'un homomorphisme diagonal unitaire
\add{sur $k$} (peut-être pas associatif), $M \to M \otimes_{k} M$, qui soit
$A$-linéaire.\marginal{(pour la structure de $A$-module sur $M \otimes M$
provenant de $A \to A \otimes A$ ; il revient au même de dire que
\struck{$A \otimes M \to M$ comp.\ avec application diagonale})} Supposons de
plus donné une « unité » i.e.\ un homomorphisme $k \to M$ qui soit un
homomorphisme de coalgèbres unitaires, i.e.\ l'image \add{$e_{M}$} de $1$
satisfait $\varepsilon_{M}(e_{M}) = 1$,
$\Delta_{M}(e_{M}) = e_{M} \otimes e_{M}$. Supposons également $M$ cocomplète,
et que l'application $u \mapsto u.e$ de $A$ dans $M$ soit injective. Alors $M$
est un $A$-module libre ; de façon plus précise, posant
$\overline{M} = M \otimes_{A} k$, il existe un isomorphisme (non canonique
…) de $A$-modules $M \xrightarrow{\sim} A \otimes_{k} \overline{M}$,
\add{tel que le diagramme}\marginal{ayant une base dont fait partie $e_{1}$}

\begin{tikzcd}
M \arrow[dr] & A \otimes_{k} \overline{M} \arrow[l, "\approx"'] \arrow[d] \\
 & \overline{M}
\end{tikzcd}

\add{commute, et que …}\note{Sa phrase s'arrête là ; le diagramme et les
mots qui l'encadrent sont de sa main.}

\emph{Exemple 1.} Supposons l'application diagonale de $A$ commutative. Alors
$A$ peut être considéré comme l'hyperalgèbre d'un groupe formel connexe
\add{$G$} sur $k$. Si $\Delta_{M}$ est associatif et commutatif, $M$ aussi est
formé des distributions à l'origine d'une variété formelle connexe \add{$V$} sur
$k$, et le fait que $A$ opère sur $M$ de façon que l'application diagonale soit
$A$-linéaire signifie que $G$ opère sur $V$. On voit donc que si
$u \mapsto u.e$ \add{de $A$ dans $M$} est injectif, i.e.\ si $g \mapsto g.a$
($a$ le « centre » de $V$) de $G$ dans $V$ est un monomorphisme, alors $M$ est
libre sur $A$. En fait, on trouvera que $G$ opère librement sur $V$ i.e.

\note{Dans la marge gauche, deux notes de sa main : « \ill{} qu'il y a un
antipodisme ? » et, encadrée, « faut-il mettre dans les hyp.\ que
$A \otimes M \to M$ est-il aussi associatif ! ».}

\page{135}

que $G \times V \to V \times V$ est un monomorphisme (c'est le lemme de
Nakayama), et on sait bien alors (SGA 3 VII$_{B}$ …) que \struck{\ill{}}
l'on a un joli quotient $W = V/G$, avec \struck{\ill{}} $V \to W$ plat et en
fait la projection structurale d'un $G_{W}$-torseur. Bien sûr la coalgèbre
$\overline{M}$ n'est autre que celle qui définit le quotient $W$ précédent. En
général, il n'est pas question de prouver que $V \simeq W \times G$ comme
variété formelle à opérateurs (i.e.\ que le torseur précédent est trivial),
i.e.\ de trouver une \struck{\ill{}} section de $V$ sur $W$, i.e.\ un relèvement
de comodules $\overline{M} \to M$ ; on peut seulement dire que tout relèvement
en tant que $k$-vectoriels définira un isomorphisme
$A \otimes_{k} \overline{M} \simeq M$ de $A$-modules. Cependant, si $G$ est
formellement lisse i.e.\ le dual de $A$ une algèbre de séries formelles
(éventuellement à un nb infini de générateurs), alors on aura splittage parfait ;
ce sera la cas en particulier en car.\ nulle.

\emph{Exemple 2.} Supposons $A$ commutative pour sa loi d'algèbre. Donc $A$ est
l'algèbre affine d'un schéma en monoïdes \add{affine} $G$. L'hypothèse sur la
filtration cocomplète est assez étrange ici, elle s'énonce en disant que pour
toute $f \in A$ (fonction sur $G$) \struck{\ill{}}, il existe un entier $n$
\add{$\geq 1$} tel que la fonction en $n$ variables $f(g_{1}g_{2}\dots g_{n})$
soit \struck{identiquement nulle} \add{somme de fonctions} à $n-1$ parmi les
variables. Si $G$ est de type fini sur $k$ \add{de car.\ $p > 0$}, je pense que
\struck{cela doit entraîner que l'unité de $G$ est point isolé, donc, comme $G$
est un groupe, que $G$ est fini sur $k$.} \struck{\ill{}} Si $G$ est commutatif
(i.e.\ $\Delta_{A}$ commutatif), alors le dual de Cartier de $A$ doit être
purement infinitésimal, i.e.\ $G$ doit être unipotent (et réciproquement). On
peut conjecturer que \struck{\ill{}} la condition en question, pour $G$ de type
fini \struck{pas néc.\ comm.}, équivaut à $G$ fini unipotent, donc si on ne
suppose plus non plus $G$ de type fini, que $G$ est \struck{profini et}
prounipotent.\marginal{(car car.\ $k > 0$ !}\marginal{unipotent}

\note{Le dernier tiers de la page, depuis « cela doit entraîner » jusqu'à
« Si $G$ est commutatif », est barré de plusieurs longues diagonales.}

\page{136}

\emph{Corollaire} Supposons que l'homomorphisme
$M \to M \otimes_{k} M$ induise un homomorphisme
$M \otimes_{(A,\Delta_{A})} (A \otimes_{k} A) \to M \otimes_{k} M$
\struck{donc par réduction} qui soit un \emph{isomorphisme}. Alors
l'homomorphisme $u \mapsto u.e$ de $A$ dans $M$ est un isomorphisme pour toutes
les structures.

Il suffit de prouver que le rang de $M$ sur $A$ est $1$. Or ce rang de
$M \otimes_{k} M$ sur $A \otimes_{k} A$ est $n^{2}$, et on aura $n = n^{2}$,
i.e.\ $n(n-1) = 0$ d'où, comme $n \neq 0$, $n = 1$. Ceci ne marche que si on
sait que le rang est fini. Mais on peut dire aussi que l'application diagonale
pour $\overline{M}$, déduite de celle de $M$ par passage au quotient (NB
$\overline{M} \otimes \overline{M} = (M \otimes_{k} M) \otimes_{(A \otimes_{k}
A)} k$) est \emph{surjective}, donc bijective puisqu'elle est unitaire,
donc tout élément de $\overline{M} \otimes \overline{M}$ est de la forme
$\overline{e} \otimes x$, ce qui n'est possible que si $\overline{e}$ engendre
$\overline{M}$.

Ceci montre que la conjecture de locale trivialité soulevée dans mes notes sur
les $\otimes$-catégories marche si la bialgèbre $A$ envisagée est finie sur $k$,
\add{et sa duale est cocomplète,} \emph{pourvu que} l'algèbre $B$ admette
un homomorphisme d'algèbres dans $k$ (ou une extension de $k$), i.e.\ pourvu que
$B$ ne soit pas identique à son idéal des commutateurs.

\section*{Motifs à coefficients dans un corps de nombres, de sa main (pages 138 à 139)}

\note{Deux feuillets entièrement de sa main, d'une autre encre et d'un autre
papier que le tapuscrit qui précède. Le titre est le sien, en tête de la page
138.}

\page{138}

\subsection*{Motifs à coefficients dans un corps de nbs.}

Soit d'abord $\underline{C}$ une $\otimes$-catégorie sur un corps $K$. Si $K'$
est une extension \add{finie} de $K$, on désignera par $\underline{C}^{K'}$ la
catégorie des objets de $\underline{C}$ sur lesquels $K'$ opère
\uncertain{linéairement}, i.e.\ des $X \in \operatorname{Ob} \underline{C}$
munis d'un homomorphisme de $K$-algèbres $K' \to \operatorname{End}(X)$. C'est
encore une $\otimes$-catégorie, en définissant $X \otimes_{K'} Y$ comme le
quotient de $X \otimes_{K} Y$ par les
$\lambda \otimes \operatorname{id}_{Y} - \operatorname{id}_{X} \otimes \lambda$
($\lambda \in K'$), et $\underline{\operatorname{Hom}}_{K'}(X,Y)$ comme le noyau
de $\underline{\operatorname{Hom}}_{K}(X,Y)$ sous
$\underline{\operatorname{Hom}}(\lambda_{X},\operatorname{id}_{Y})
- \underline{\operatorname{Hom}}(\operatorname{id}_{X},\lambda_{Y})$.

Quand $\underline{C}$ est défini par un schéma en groupes affine $G$ sur $K$
\add{muni un foncteur fibre $F$ sur $\underline{C}$, \ill{}}, alors
$\underline{C}^{K'}$ est défini par le \struck{groupe} \struck{\ill{}}
$G_{K'} = G \otimes_{K} K'$,\marginal{sur $\underline{C}^{K'}$} qui a un
foncteur fibre $F_{K'}$ déduit trivialement de $F$.

On a un $\otimes$-foncteur naturel
$\underline{C} \to \underline{C}^{K'}$, notant $X \mapsto X \otimes_{K} K'$, et
dont le composé avec le foncteur « oubli » est bien
$X \otimes_{K} \widetilde{K'}$, où $\widetilde{K'}$ est \struck{\ill{}}
l'espace vectoriel \emph{sous-jacent} à $K'$ sur $K$, --- puisque
\ill{} $\underline{C}$ défini par le \ill{}
$\operatorname{Vect}$ \add{de dim.\ finie} sur $K \to \underline{C}$ ]. \ill{}
est un $\otimes$-foncteur évident \ill{}

Supposons $\underline{C}$ la catégorie des motifs \add{(sur)}
\uncertain{semi-simples} (\ill{}) sur un corps $k$, donc
$K = \mathbb{Q}$, $K'$ un corps de nbs. Soit $E$ un motif à coefficients dans
$K'$, donc \ill{} $K' \to \operatorname{End}_{\mathbb{Q}}(E_{0})$, $E_{0}$ le
$\mathbb{Q}$-motif \ill{}, avec \marginal{un motif dans $\underline{C}^{K'}$}
\ill{}. Alors $\widetilde{E}$ \add{est} défini par $\widetilde{E}$ dans
$\underline{C}$, avec l'hom\add{omorphisme}

\page{139}

donné par
\[ K' \xrightarrow{\ \varphi\ } \operatorname{End}_{\mathbb{Q}}(E_{0})
   \xrightarrow{\ \ill{}\ }
   \operatorname{End}_{\mathbb{Q}}(\check{E}_{0}) . \]
\note{La flèche de droite porte au-dessus d'elle une glose interlinéaire de
trois ou quatre lettres qui ne se laisse pas lire.}
Donnons-nous une polarisation sur $E_{0}$ (ou \ill{} \struck{\ill{}} : $E$).
Alors on trouve un isom.\ $\check{E}_{0} \simeq E_{0}(\rho)$ ($E_{0}$ de poids
$\rho$), donc
\struck{$K' \to \operatorname{End}_{\mathbb{Q}}(\check{E}_{0})$}
\[ \operatorname{End}_{\mathbb{Q}}(\check{E}_{0})
   \simeq \operatorname{End}(E_{0}(\rho)) \simeq \operatorname{End}(E_{0}) , \]
et \ill{} \ill{} une \emph{involution} \add{\ill{}} sur
$\operatorname{End}(\check{E}_{0})$. Le motif dual $(E_{\varphi})^{\vee}$ de
$E_{\varphi}$ (défini par $\varphi$) est défini par l'hom.\
$\lambda \mapsto \varphi(\lambda)' \uncertain{=} \varphi'(\bar{\lambda})$,
\ill{}
\[ \boxed{\ (E_{\varphi})^{\vee} = (E_{\varphi'})(\rho) \ } \]
\note{L'égalité est encadrée sur le feuillet.}

Ce n'est pas en général un isom.\
$(E_{\varphi})^{\vee} \simeq E_{\varphi}(\rho)$ !

Prenons l'exemple des $K'$-motifs de niveau zéro, poids $0$. La catégorie est
équivalente à la catégorie des \struck{\ill{}} vectoriels de dim.\ finie sur
$K'$, \struck{\ill{}} \ill{} à opération \add{continue} de
$\pi = \operatorname{Gal}(\bar{k}/k)$. Le dual correspond à l'opération
contragrédiente de $\pi$. Les $K'$-motifs de rang $1$ \add{à isom.\ près}
correspondent aux représentations \add{continues} \struck{de}
$\pi \xrightarrow{\ \chi\ } K'^{*}$, et on a
\[ (L_{\chi})^{\vee} = L_{\check{\chi}} , \quad \text{où} \quad
   \check{\chi}(g) = \chi(g^{-1}) . \]
Donc on a. $L_{\chi}^{\vee} \simeq L_{\chi}$ sss $\chi = \check{\chi}$ i.e.\
$\chi$ à valeurs dans $\pm 1$ !

\section*{Structures supplémentaires sur des vectoriels, et algèbres de Hopf (page 140)}

\note{Chemise de papier brun, donnée ici sans numéro de page : elle porte, de sa
main, le titre du faisceau qui commence après ce lot. Une première ligne,
au-dessus, est raturée jusqu'à l'illisible ; elle commençait par « Fasc… ».}

\struck{\ill{}}

Structures supplémentaires sur des vectoriels, et algèbres de Hopf
\uncertain{top.}

Anneaux des représentations de certains groupes algébriques.

\end{document}
